\documentclass[11pt]{article}
\usepackage[margin=1.1in]{geometry}
\usepackage{amsmath,amssymb,amsthm,booktabs,longtable,xcolor,array}
\usepackage{listings}
\usepackage[colorlinks=true,linkcolor=blue!50!black,citecolor=blue!50!black,urlcolor=blue!50!black]{hyperref}
\usepackage{orcidlink}

\newtheorem{fact}{Fact}
\newtheorem{prop}{Proposition}
\newtheorem{conj}{Conjecture}
\newcommand{\lam}{\lambda}
\newcommand{\code}[1]{\texttt{#1}}
\emergencystretch=2em
\hyphenation{re-flec-tion-ad-mit-ting quasi-crys-tal-line}

\lstdefinelanguage{lean}{
    morekeywords={theorem,def,structure,instance,namespace,end,deriving,
                  by,decide,omega,simp,induction,intro,Int,Nat,List,Prop},
    morecomment=[l]{--},
    morecomment=[n]{/-}{-/},
    sensitive=true
}
\lstset{basicstyle=\ttfamily\small,breaklines=true,frame=single,
        framerule=0.3pt,columns=fullflexible,keepspaces=true,
        showstringspaces=false,commentstyle=\color{gray}\itshape,
        keywordstyle=\color{blue!60!black}\bfseries}

\title{\vspace{-1cm}
\textbf{Area and Perimeter Inflation in the Spectre Monotile}}

\author{Brent S. Hartshorn \orcidlink{0009-0004-2853-655X}
        (\url{brenthartshorn@proton.me})}

\begin{document}
\maketitle


\begin{abstract}
The Spectre substitution on nine species is studied as an exact object.  Its
transfer spectrum is $x^{5} (x - 1)(x + 1)(x^{2} - 8 x + 1)$, so the Perron eigenvalue is
$\lam^{2}=\sqrt{15} + 4$, the fundamental unit of $\mathbb{Z}[\sqrt{15}]$, and the
linear inflation factor $\lam=\tfrac12(\sqrt6+\sqrt{10})$ splits into a
hexagonal and a pentagonal channel without either square root being evaluated.
The zero sector is not semisimple: rank $M=5$ but rank $M^{n}=4$ for
$n\ge2$, so one transient direction survives a single substitution and dies
only after two.  Two marginal eigenvalues carry exact integer charges, one
conserved and one alternating, and the alternating one tracks global handedness
from species counts alone.  The census is stated twice, as a matrix count and
as a placed-tile count, and the relation between them is derived rather than
stipulated: row $\Gamma$ is all ones, so the Mystic contributes exactly the
previous generation and $T_{\rm geom}(n)=T(n)+T(n-1)$.  We give the tile area
as an exact quadratic form, the Spectre-to-Hat order parameter in closed form
in $\mathbb{Q}(\sqrt{15})$, and the Seifert minors of $T(2,k)$ in closed form
for every $k$.  Separately from the area grading we derive the operator
governing the supertile \emph{perimeter}: an exact order-three recurrence whose
companion matrix has characteristic polynomial $(x - 1)(x^{2} - 4 x - 1)$, so the boundary grows by
$\nu=2+\sqrt5=\varphi^{3}$ while the area grows by $\lam^{2}=4+\sqrt{15}$ ---
the same $\mathbb{Q}(\sqrt{15})$ / $\mathbb{Q}(\sqrt5)$ divide, now between
area and perimeter of one tile.  The boundary is therefore fractal, of
dimension $1.39925321391$, and the resemblance to the Hat's $\varphi^{4}$ is shown to be a
Lucas identity rather than a shared operator.  The perimeter sequence is
counted twice, the second time in exact $\mathbb{Q}(\sqrt3)$ arithmetic, so no
result here rests on a floating-point tolerance.  
\end{abstract}


\section{Introduction}
\label{sec:intro}

The motivation is a question about the initial state of the universe.  Work on the
Nariai horizon, the de~Sitter spacetime $dS_2\times S^2$ whose bifurcation
spheres are compact, has related the Araki--Uhlmann relative entropy of a
coherent excitation to the modular-weighted boost-energy flux through the
horizon, and the $s$-wave sector of the horizon dynamics, linearized
Jackiw--Teitelboim gravity, is consistent with the area law
$S_{\rm rel}=\delta A/4G$.  That is the converse direction of the result of
Dorau and Much~\cite{DorauMuch}: the Einstein equations imply the area law.  It
does not derive the normalization, whose coefficient is the Newton constant,
and the sign of $\delta A$ depends on the boundary conditions.  Nothing in the
present paper depends on that result; it is the setting in which the question
arises.  The thesis here is that the crossover surface --- a two-dimensional
conformal carrier with vanishing configurational entropy density --- is a
substitution monotile phase \cite{SmithEtAl, SmithEtAl2}.  
That identification is a conjecture and is labelled as one.  What can
be done without it, and what occupies most of what follows, is to compute the
substitution's exact structure: the identification is only worth making if the
object identified has content, and the content is arithmetic.

Statements are stratified by epistemic weight.  \emph{Facts} are computations
or theorems.  \emph{Propositions} are computations with a modelling step.
\emph{Conjectures} are sharp claims not yet proved.  A reader who rejects every
conjecture here still keeps every Fact, and the Facts are the reason to read on.
The Facts that are closed statements are not merely computed twice.
They are proved.  \code{spectrelean.py} emits them as Lean~4 and the kernel
checks them, which is a different kind of evidence from a second computation
and is reported separately in Section~\ref{sec:lean}.


\section{What this is a substrate for}
\label{sec:bigpicture}

The results below are arithmetic: two algebraic units, the fields they
generate, and the operators whose spectra they are.  A reader coming from
physics is owed an account of why arithmetic of this kind is the right thing
to compute, and where each piece of it is consumed.  This section is that
account.  It contains no new Facts; it says what the Facts are for.

\subsection*{Why a tiling carries a unit at all}

A substitution tiling is a rule for replacing every tile by a patch of tiles
similar to it, applied without end.  Counting tiles under the rule is a linear
map, its matrix has a leading eigenvalue, and by Perron--Frobenius that
eigenvalue is an algebraic integer.  For the Spectre it is
$\lam^{2}=\sqrt{15} + 4$, the fundamental unit of $\mathbb{Z}[\sqrt{15}]$
(Section~\ref{sec:spectral}).  Nothing about that number was chosen: it is
forced by the rule, and the rule is forced by the tile.

The unit is Pisot --- its conjugate $4 - \sqrt{15}$ lies inside the unit circle --- and
the Pisot property is what turns geometry into arithmetic.  Bombieri and
Taylor, and Solomyak after them, showed that a Pisot inflation gives a tiling
whose diffraction is pure point: it scatters like a crystal despite having no
period.  Bellissard's gap-labelling theorem then says that the integrated
density of states of any Hamiltonian on the tiling takes its values in the
frequency module, which for the Spectre is $\mathbb{Z}[1/\lam^{2}]=
\mathbb{Z}[\sqrt{15}]$.  So the field is not decoration.  It is the set of
numbers an experiment on this tiling is permitted to return, whether that
experiment is a diffraction pattern or a spectral gap.

\subsection*{Two fields, and what separates them}

The Hat, the other aperiodic monotile, inflates in $\mathbb{Q}(\sqrt5)$
(Section~\ref{sec:phases}).  A diffraction measurement can tell the two fields
apart, and the 2026 measurement cited there does.  That is the sense in which
the tiling is a \emph{phase}: two monotiles, two fields, and an observable
that reads which one is present.

What Section~\ref{sec:bdyop} adds is that the second field was already inside
the first tile.  The Spectre's area inflates in $\mathbb{Q}(\sqrt{15})$; its
perimeter inflates by $\nu=2+\sqrt5=\varphi^{3}$, in $\mathbb{Q}(\sqrt5)$.  One
tile, two units.  The consequence for physics is a rule of thumb that the
rest of this paper makes exact: \emph{the field an observable lands in tells
you what kind of quantity it is.}  Counts, densities, areas, integrated
densities of states --- anything that scales with the number of tiles --- are
$\mathbb{Q}(\sqrt{15})$ quantities.  Lengths, interfaces, contact lines,
anything that scales with a boundary --- are $\mathbb{Q}(\sqrt5)$
quantities.  Their ratio, which a multi-scale argument needs, lives in
neither and is a degree-four number in the compositum
$\mathbb{Q}(\sqrt3,\sqrt5)$ (Proposition~\ref{prop:compositum}).

The boundary is also fractal, of dimension $1.39925321391$, because $\nu$ exceeds the
linear inflation $\lam$.  A supertile's outline grows faster than its size.
That is the geometric content of ``two units'': a shape whose perimeter and
area do not scale together cannot be assigned a single inflation factor, and
the mismatch is measured by the ratio of the two fields' units.

A fractal with a Hausdorff dimension of $1.39925321391$ represents a geometric structure with a complexity that sits precisely between a 1D line and a 2D plane.
In nature and medicine, fractal dimensions around 1.40 are frequently observed in the highly folded structural boundaries of the human brain (such as the boundary between the cerebral cortex and white matter \cite{bullmore}).

\subsection*{Where the substrate is consumed}

Three programmes draw on this.  The first is the horizon setting of the
introduction~\cite{DorauMuch}.  There the Nariai throat's
crossover surface is a two-dimensional carrier with zero configurational
entropy density, and the conjecture --- labelled as one in
Section~\ref{sec:data} --- is that it is a substitution monotile phase \cite{SmithEtAl, SmithEtAl2}.  
If so, every exact quantity here is an exact quantity of the horizon: its
gap labels, its charges, its two units.  The identification is a modelling
step; the arithmetic is not.

The second is spectral.  The integer charges of Section~\ref{sec:charges},
the X-charge of Section~\ref{sec:mass}, and the gap-label module are all
statements about what a Hamiltonian on the tiling can and cannot do, and they
are exact, so a numerical spectrum that disagrees with them is wrong rather
than approximate.  That is a useful kind of constraint to have.

The third is the multi-scale programme of Section~\ref{sec:kakeya}, and it
is recorded here as a negative result because that is what it is.  The
Kakeya machinery of Wang and Zahl wants a direction set that is spread;
aperiodicity supplies non-repetition, which is a different property, and a
periodic lattice beats the monotile on the quantity Kakeya actually needs.
What the substitution does have is scales, graded by the unit, and along that
ladder the monotile behaves as a sticky family --- which in the Kakeya
programme is the obstruction a proof must survive, not a help.  A reader
hoping the aperiodic substrate would make Kakeya easier should stop hoping
here; a reader wanting a concrete, finite, computable sticky family now has
one.

\subsection*{How to read the rest}

The Facts are exact and most are proved in Lean.  The Propositions carry one
modelling step each, named where it occurs.  The Conjectures are sharp and
unproved.  The negative results are not softened.  Every number in this document is
generated from the same code that checks it, so the text and the arithmetic
cannot disagree.


\section{The crossover screen and the two monotile phases}
\label{sec:screen}

A zero-entropy two-dimensional carrier admitting quasi-crystalline long-range
order must, if built from a single tile, be one of the known aperiodic monotile
phases, and the two available phases differ in exactly one theorem-level
respect: the Spectre phase is strictly chiral --- reflections are globally
forbidden --- while the Hat phase requires both handednesses.  An isolated
two-dimensional screen has no ambient notion of reflection: a chirality, once
fixed by the matching rules, is absolute.  A leaf embedded in a
three-dimensional bulk does, since reflection of a planar patch is an ambient
rotation.

\begin{conj}[Chirality relaxation]
\label{conj:chirality}
The admissible monotile phase on an isolated two-dimensional conformal screen
is the strictly chiral Spectre phase; the admissible phase on a leaf embedded
in an emergent three-dimensional bulk is the reflection-admitting Hat phase.
The Spectre-to-Hat transition is a geometric order parameter for the emergence
of the third spatial dimension: parity, forbidden on the screen, is restored by
embedding.
\end{conj}

This is a selection statement, not a dynamics, and the selection criterion ---
can reflections be realized ambiently? --- is exactly the criterion by which
the two phases differ, and by nothing else.  Everything that follows equips the
conjecture with exact structure rather than defending it: an integer census
charge locked to handedness (Section~\ref{sec:charges}), an order parameter
along the Tile$(a,b)$ continuum (Section~\ref{sec:deform}), 
and a second, arithmetic distinction between the phases
(Section~\ref{sec:phases}) that the parity argument does not supply.


\section{Exact spectral theory of the substitution}
\label{sec:spectral}

The substitution acts on nine species $(\Gamma,\,\Delta,\,\Theta,\,\Lambda,\,\Xi,\,\Pi,\,\Sigma,\,\Phi,\,\Psi)$, with $\Gamma$ carrying the Mystic pair
$\Gamma_1+\Gamma_2$, $\Gamma_2=\mathrm{Tile}(b,a)$.  The matrix is the one
quoted input of this paper; everything else is derived from it.
{\setlength{\arraycolsep}{2.6pt}
\begin{equation}
M=\begin{pmatrix}1&1&1&1&1&1&1&1&1 \\1&1&1&1&1&1&1&1&1 \\1&0&0&0&0&0&0&0&0 \\0&0&0&0&0&0&1&0&0 \\1&2&0&1&0&1&2&0&0 \\1&1&2&1&1&0&1&1&0 \\1&1&1&1&1&1&1&1&1 \\1&2&2&2&2&2&1&2&2 \\0&0&1&1&2&2&0&2&3\end{pmatrix}
\label{eq:M}
\end{equation}}

\begin{fact}[Inflation factor and its arithmetic]
\label{fact:lambda}
The Perron eigenvalue is $\lam^{2}=\sqrt{15} + 4$, the fundamental unit of
$\mathbb{Z}[\sqrt{15}]$: $4^{2}-15\cdot1^{2}=1$ is the least Pell solution.
The linear inflation factor is
\[
\lam=\sqrt{\sqrt{15} + 4}=\tfrac12(\sqrt6+\sqrt{10})=2.8058837\ldots,
\qquad
(\lam-\lam^{-1})^{2}=6,
\qquad
(\lam+\lam^{-1})^{2}=10,
\]
with minimal polynomial $x^{4} - 8 x^{2} + 1$ and splitting field $\mathbb{Q}(\sqrt6,\sqrt{10})$.
The two channel identities are proved in $\mathbb{Z}[\sqrt{15}]$ without
evaluating a square root: each is $\lam^{2}+\lam^{-2}\mp2$, and the
$\sqrt{15}$ parts cancel.  Inflation therefore acts on all area-graded data by
multiplication by a fundamental unit.
\end{fact}

\begin{fact}[Transfer spectrum: a Jordan pair among the zero modes]
\label{fact:spectrum}
$\chi_M(x)=x^{5} (x - 1)(x + 1)(x^{2} - 8 x + 1)$, and the zero sector carries structure the spectrum alone does
not show:
\[
\operatorname{rank}M=5,\qquad \operatorname{rank}M^{n}=4\quad(n\ge2).
\]
The zero eigenvalue has algebraic multiplicity five and geometric multiplicity
four, so $M$ is not diagonalizable.  The transient sector is four census
directions erased in one substitution together with one Jordan pair erased in
exactly two: a one-step afterimage, the only place where the substitution
remembers longer than its eigenvalues admit.  The remaining spectrum is the
Pisot-unit pair $\{\lam^{2},\lam^{-2}\}$ and two marginal eigenvalues $\pm1$.
Deviations from equilibrium contract by $\lam^{-4}=(4 - \sqrt{15})^{2}\approx0.0161332$ per step,
except along the marginal directions, which are exact integer charges, and
along the Jordan pair, which is gone after two.
\end{fact}


\section{The census, twice, and the identity between them}
\label{sec:census}

A patch can be counted two ways and the difference matters, so it is settled
here before anything uses it.

The \emph{matrix census} is what iterating~\eqref{eq:M} produces: nine species,
one $\Gamma$ entry.  The \emph{geometric census} is what a placed patch
contains, and a placed patch carries the Mystic as well, so it is larger.  The
previous draft quoted the second while computing with the first.

\begin{fact}[The Mystic is the previous generation]
\label{fact:mystic}
Row $\Gamma$ of $M$ is all ones, so $N_\Gamma(n)$ is the entire census at depth
$n-1$; the Mystic is a second copy of it.  Hence
\[
T_{\rm geom}(n)=T(n)+T(n-1),
\]
and both sequences obey $T(n{+}1)=8T(n)-T(n{-}1)$, the recurrence whose
characteristic polynomial is the quadratic factor of $\chi_M$.
\end{fact}

\begin{table}[h]
\centering
\begin{tabular}{rrrr}
\toprule
depth & matrix census & geometric census & Mystic \\
\midrule
0 & 1 & 1 & 0 \\
1 & 8 & 9 & 1 \\
2 & 63 & 71 & 8 \\
3 & 496 & 559 & 63 \\
4 & 3905 & 4401 & 496 \\
\bottomrule
\end{tabular}
\caption{The two censuses. The fourth column is the third minus the second,
and is also the second column one row up --- which is Fact~\ref{fact:mystic}.
Charges are linear forms on the nine-species vector and are evaluated against
the matrix column.}
\end{table}


\section{Two integer charges}
\label{sec:charges}

\begin{fact}[The marginal charges]
\label{fact:charges}
The left eigenvectors for the marginal eigenvalues $\pm1$, in the species order
of~\eqref{eq:M}, are
\[
Q_{+}=(-3, -1, 0, 0, 1, 1, -2, 1, 2),\qquad Q_{-}=(-1, -1, 2, 0, 1, -1, 0, 1, 0),
\]
satisfying $Q_{+}M=Q_{+}$ and $Q_{-}M=-Q_{-}$ exactly over $\mathbb{Z}$.
Evaluated on the matrix census of Section~\ref{sec:census} at depths one to
four,
\[
Q_{+}=-1, -1, -1, -1,\qquad Q_{-}=+1, -1, +1, -1.
\]
$Q_{+}$ is conserved; $Q_{-}$ alternates.
\end{fact}

\begin{prop}[$Q_{-}$ is the parity's census avatar]
\label{prop:parity}
The handedness census of the placed tiles gives a chirality order parameter
$\chi=\pm1$ at every depth --- the tiling is strictly single-handed at every
scale --- and the sign of $Q_{-}$ alternates in lockstep with it.  Global
geometric parity is computable from species counts alone.  Transcribing the
published placement chain into $\mathbb{Q}(\sqrt3)$ affine arithmetic, every
determinant is exactly $\pm1$ and the common determinant of a depth-$n$ patch
equals $-Q_{-}(n)$ at every depth computed.
\end{prop}

\begin{fact}[The generational ledger]
\label{fact:ledger}
Row $\Theta$ of $M$ is the indicator of $\Gamma$ and row $\Lambda$ the
indicator of $\Sigma$, so
\[
N_\Theta(n)=N_\Gamma(n-1),\qquad N_\Lambda(n)=N_\Sigma(n-1)
\]
hold exactly, not asymptotically.  Two species are a verbatim copy of an
earlier generation, which gives the modular arrow an integer bookkeeping.
\end{fact}


\section{Deformations and the order parameter}
\label{sec:deform}

\begin{fact}[The area form, and the order parameter in closed form]
\label{fact:area}
The area of $\mathrm{Tile}(a,b)$ is the exact quadratic form
\[
A(a,b)=2 \sqrt{3} a^{2} + 3 a b + \sqrt{3} b^{2},
\]
an identity in $\mathbb{Q}(\sqrt3)[a,b]$ obtained by symbolic shoelace on the
published vertex construction, not fitted to sample points.  It reproduces the
endpoints $A(1,\sqrt3)=8 \sqrt{3}$ and $A(\sqrt3,1)=10 \sqrt{3}$, and gives the Mystic asymmetry
exactly:
\[
A(b,a)-A(a,b)=\sqrt{3} (- a^{2} + b^{2}),
\]
which is the precise content of ``only $\Gamma$ moves''.  With the Mystic pair
carrying both areas, the splitting is maximal at the Hat endpoint with exact
value
\[
\max_{r\in[1,\sqrt3]}\bigl|f_\Gamma-f_\Delta\bigr|
=\frac{10g}{8+10g}=\frac{35}{67} - \frac{20 \sqrt{15}}{201}=0.13701658\ldots,
\qquad g=4 - \sqrt{15}.
\]
The $\sqrt3$ of the geometry cancels and the maximum lies in
$\mathbb{Q}(\sqrt{15})$ alone: the hexagonal world supplies the areas, the
Pisot world the frequencies, and the order parameter forgets the former.
\end{fact}


\section{Two phases, two quadratic fields}
\label{sec:phases}

Conjecture~\ref{conj:chirality} distinguishes the two monotile phases by
parity: the Spectre phase is strictly chiral, the Hat phase admits both
handednesses.  There is a second distinction, arithmetic rather than
geometric, and it has not been used before.

\begin{fact}[The phases are graded by different fields]
\label{fact:phases}
The Hat is generated by a substitution on four metatiles $H,T,P,F$, not nine
\cite{SmithEtAl}.  Its substitution matrix, in the order $H,T,P,F$ and up to
transposition (which does not change the characteristic polynomial), is
$\left(\begin{smallmatrix}3&1&2&2\\1&0&0&0\\3&0&1&1\\3&0&2&3\end{smallmatrix}\right)$,
and it gives
\[
\chi_H(x)=(x - 1)(x + 1)(x^{2} - 7 x + 1),
\qquad
\kappa=\tfrac12(7+3\sqrt5)=\varphi^{4}=6.85410196625\ldots,
\]
against the Spectre's $\chi_M(x)=x^{5} (x - 1)(x + 1)(x^{2} - 8 x + 1)$ and $\lam^{2}=\sqrt{15} + 4=7.87298334621\ldots$.  The two
characteristic polynomials share the factor $(x-1)(x+1)$ and differ in a single
coefficient of the quadratic --- seven against eight --- and that coefficient
is the discriminant.  The Spectre is graded by $\mathbb{Q}(\sqrt{15})$ and the
Hat by $\mathbb{Q}(\sqrt5)$.
\end{fact}

\begin{fact}[Trace accounting, in two fields]
\label{fact:trace}
In both phases the marginal pair $\pm1$ cancels and the zero modes contribute
nothing, so the trace is the Perron unit plus its conjugate:
\[
\operatorname{tr}M=\lam^{2}+\lam^{-2}=8,
\qquad
\operatorname{tr}H=\kappa+\kappa^{-1}=7.
\]
The same accounting, run in $\mathbb{Q}(\sqrt{15})$ and in
$\mathbb{Q}(\sqrt5)$.
\end{fact}

\subsection{Aperiodicity as irrationality, and which tile witnesses it}

A tiling with a period is a periodic arrangement of a fundamental domain; in a
fundamental domain every species occurs a whole number of times; so in a
periodic tiling every relative frequency is rational.  One irrational frequency
therefore refutes periodicity.  The two phases differ in how much of the tiling
can do the refuting.

\begin{fact}[Every Spectre species is a witness; one Hat metatile is not]
\label{fact:witness}
All 9 Spectre frequencies lie in $\mathbb{Z}[\sqrt{15}]$ with nonzero
$\sqrt{15}$ part, so any one of them refutes periodicity.  The Hat's four
metatile frequencies are
\[
f_H=\frac{1}{3},\qquad f_T=\frac{7}{6} - \frac{\sqrt{5}}{2},\qquad f_P=-2 + \sqrt{5},\qquad f_F=\frac{3}{2} - \frac{\sqrt{5}}{2},
\]
and the first is rational.  The commonest Hat metatile witnesses nothing, and
the argument has to be run on a rarer one.
\end{fact}

\noindent
That the Spectre frequencies are algebraic \emph{integers} rather than merely
algebraic numbers is what lets the eigenvector equation be checked in
$\mathbb{Z}[\sqrt{15}]$ with no rational arithmetic at all, which is how
Section~\ref{sec:lean} proves it rather than recomputing it.

\subsection{Measured, in a laboratory}

The Hat phase is no longer only a mathematical object.  Moritake
\emph{et al.}~\cite{Moritake2026} fabricated a monotile quasilattice --- the
centroids of the $T$ metatiles of an $H_6$ patch, 372{,}100 holes etched in
silicon nitride over half a millimetre --- and measured its optical
diffraction.  The pattern shows sharp Bragg peaks whose positions do not move
with the illumination spot, which is the experimental signature of
quasicrystalline order, and it is \emph{chiral}: a pinwheel with no mirror
line, reversing handedness when the structure is mirrored, and responding
differently to left- and right-handed circular polarization.  That last
response is absent in every mirror-symmetric quasicrystal, the Penrose tiling
included.

For the present paper the useful part is that they give the twist in closed
form.  The metatiles rotate away from the underlying honeycomb by an angle
whose coefficients follow the Fibonacci recurrence, so the ratio tends to
$\varphi^{2}$ and the limiting twist is
\[
\theta_{\rm chiral}=\arccos\frac{3\varphi-1}{4}
=\arccos\frac{1+3\sqrt5}{8}=15.5225\ldots^{\circ}.
\]

\begin{fact}[The measured angle is the inflation factor]
\label{fact:chiral}
$3\varphi-1$ and $\varphi^{4}-3$ are the same number, so
\[
\cos\theta_{\rm chiral}=\frac{\kappa-3}{4},
\qquad \kappa=\varphi^{4},
\]
with $\kappa$ the Hat inflation factor of Fact~\ref{fact:phases}.  The angle a
laboratory measured in a diffraction pattern and the Perron eigenvalue of the
substitution matrix are one constant, shifted and scaled.
\end{fact}

\noindent
This matters for Conjecture~\ref{conj:chirality} in a way the conjecture did
not anticipate.  The Hat phase's chirality is now an observed
laboratory property with an exactly predicted magnitude, and that magnitude is
fixed by the arithmetic of $\mathbb{Q}(\sqrt5)$.

It also yields a prediction rather than a restatement.  Substituting the
Spectre's inflation factor into the same formula gives
$(\lambda^{2}-3)/4=1.21825>1$, which is not a cosine: the Spectre cannot share the
Hat's twist relation, and its own chiral angle --- whatever it is --- belongs
to $\mathbb{Q}(\sqrt{15})$ instead.  We do not compute it here, because doing
so needs the Spectre's metatile twist geometry and not merely its eigenvalue.
Fabricating the Spectre quasilattice and measuring a \emph{different} chiral
angle would be a direct test of the two-field claim of
Fact~\ref{fact:phases}, at optical bench scale, with no cosmology involved.

This sharpens the conjecture at no cost.  If the Spectre-to-Hat transition is
the emergence of the third spatial dimension, then it is not only a parity
restoration but a change of the quadratic field that grades all area data ---
and the log-periodic fingerprint changes with it, from period
$\log(4+\sqrt{15})$ on the screen to $\log\varphi^{4}$ in the bulk.  Two
frequencies, named in advance, rather than one.  Whether either is present is a
question for data; that they differ is a Fact.


\section{The knot layer, for every $k$ at once}
\label{sec:knots}

\begin{fact}[Seifert minors in closed form]
\label{fact:sandwich}
The Seifert form of $T(2,k)$ is the $(k-1)$-dimensional tridiagonal
$(-2,1,1)$.  Its leading principal minors obey $D_{m}=-2D_{m-1}-D_{m-2}$ and
therefore
\[
D_m=(-1)^{m}(m+1),
\]
so they alternate in sign and never vanish.  By Sylvester's criterion the form
is negative definite at every size, giving $\sigma(T(2,k))=-(k-1)$ for every
odd $k$, hence $u\ge(k-1)/2$; one crossing change sends $\sigma^{k}$ to
$\sigma^{k-2}$, and induction closes the sandwich at
$u(T(2,k))=\tfrac{k-1}{2}$.
\end{fact}

\noindent
The closed form is proved by induction in Section~\ref{sec:lean}, which is
what upgrades this from a check at sampled $k$ to a statement about the
family.  Additivity of the signature under connected sum then makes the
selection rule exact arithmetic: $\sigma(K\#K)=-2(k-1)$ forces
$u(K\#K)=k-1$, so same-handed pairs cannot bind, while $\sigma(K\#\overline K)=0$
and the lower bound is lost.  Binding is a mirror-channel phenomenon as a
matter of signatures --- and the strictly chiral screen, having no mirror
channel, cannot bind at all.


\section{The X-charge, and the mirror channel}
\label{sec:mass}

Everything so far has been arithmetic on a nine by nine matrix.  The X-charge
is not: it is a statement about which edge of a placed tile meets which edge of
its neighbour, and no eigenvector recovers it.  The tiling has to be built.

Each tile is a fourteen-gon whose slots carry canonical roles in the cyclic
sequence \code{aabbaabbaaaabb}.  A shared edge is an \emph{X-bond} when the two
tiles disagree about the role of the slot they are gluing.  This is a statement
about roles and not about lengths: a tiling cannot glue edges of different
length, so typing the bond by realized length makes every bond match by
construction and the charge collapses to zero everywhere.  The two readings
differ on exactly one species, and that difference is the content of what
follows.

\begin{fact}[The contact atlas stabilizes]
\label{fact:atlas}
The unordered adjacency classes --- which species meets which, at which pair of
slots --- number 131 at depths 3, 4, 5, 6, with none gained and none lost.  Per-slot
statements are therefore statements about the infinite tiling and not about a
patch.
\end{fact}

\begin{fact}[X-charges, and their parity]
\label{fact:xcharge}
On interior tiles the X-status of every slot is deterministic for every species
but one:
\begin{align*}
&\underbrace{X_{\Gamma_2}=14}_{\text{Mystic, saturated}}
&&\underbrace{X_{\Gamma_1}=X_\Delta=X_\Sigma=4}_{\text{triplet}}
&&\underbrace{X_\Lambda=2}_{\text{splits }\{\Theta,\Lambda\}}\\[1ex]
&\underbrace{X_\Phi\in\{0,2\}}_{\text{flavour doubling}}
&&\underbrace{X_\Theta=X_\Pi=X_\Xi=X_\Psi=0}_{\text{sterile quartet}}
\end{align*}
Identical at every depth computed, and every value even, so the binding
spectrum is $2\delta\{0,1,2,7\}$.  Mystic saturation has a one-line reason:
$\mathrm{Tile}(b,a)$ realizes every canonical role at the other length, so
every bond it forms is a role mismatch.
\end{fact}

\begin{fact}[The flavour frequency is not converged]
\label{fact:pdrift}
$\Phi$'s two mixed slots are perfectly correlated --- never one --- so $\Phi$
splits into flavours $\Phi^{0},\Phi^{2}$.  The collared frequency
$p=\Pr(\Phi^{2})$ measured at depths 3, 4, 5, 6 is
\[
p=0.683,\ 0.579,\ 0.538,\ 0.519,
\]
still falling.  A value near $0.58$ is what depth four gives; it is a
measurement at a depth, not a limit, and earlier drafts quoted it as though it
were one.
\end{fact}

\begin{conj}[The bulk X-density in closed form]
\label{conj:rho}
If $p\to\tfrac12$ then
\[
\rho_X=\frac{2g}{1+g}=1-\frac{\sqrt{15}}{5}=0.22540333\ldots,
\qquad g=4 - \sqrt{15},
\]
an element of $\mathbb{Q}(\sqrt{15})$ with no free parameter.  Repeated Aitken
extrapolation of the measured values gives $0.5107$ then $0.5035$, which is
evidence for one half and not a proof of it.  This is the open problem of the
previous draft, reduced from ``compute $p$'' to ``decide whether $p=\tfrac12$''.
\end{conj}

\subsection{Why the mechanism switches on where it does}

Brittenham and Hermiller~\cite{BH} proved that unknotting number is not
additive: for $K=7_1=T(2,7)$ with $u(K)=3$,
$u(K\#\overline K)\le5<6$.  Three things about that theorem bear on the
mechanism, and all three are easy to overstate.

\emph{The deficit is bounded, not determined.}  Theirs is an upper bound;
Scharlemann \cite{Scharlemann} gives $u\ge2$; so $\delta=6-u$ lies in $[1,4]$ and its value is
their own Question~4.4, open.  The binding spectrum is parameter-free in its
\emph{ratios} $\{0,1,2,7\}$ --- the X-charges halved --- and says nothing
about the scale.

\emph{The threshold is theirs.}  Their Corollary~1.3 covers $T(2,2k+1)$ for
$k\ge3$, index seven up; $T(2,3)$ and $T(2,5)$ are outside it and whether
either admits any partner is open.  So $k=7$ is where non-additivity is known
to begin, and it is also where the substrate's edge braids close.  That
coincidence is the whole of the mechanism's appeal, and it is a coincidence
until someone explains it.

\emph{The proof was corrected, and the correction was about mirroring.}  The
first version routed through two diagrams of $K15n81556$ asserted to be one
knot; Wang and Zhang~\cite{WZ} showed by the Jones polynomial that they are a
chiral knot and its mirror, and verified the inequality directly anyway.  The
current version supplies the missing isotopy and credits them.  The theorem is
unaffected and now stands by two routes, and the deficit bound does not move.
The episode is a caution, not a result: it is \emph{not} evidence for a
mirror-sensitive mechanism, and no proof layer here would have caught it,
because the question is about knot diagrams and this kernel sees arithmetic.

Additivity of the signature then makes the selection rule arithmetic:
$\sigma(K\#K)=-2(k-1)$ forces $u(K\#K)=k-1$, so same-handed pairs cannot bind,
while $\sigma(K\#\overline K)=0$ and the bound is lost.  Binding is a
mirror-channel phenomenon as a matter of signatures, and the strictly chiral
screen has no mirror channel.


\section{The other unit: a boundary operator in $\mathbb{Q}(\sqrt5)$}
\label{sec:bdyop}

Everything above grades the substitution by area.  The perimeter is graded by
something else, and the two do not live in the same field.

Let $P(k)$ be the number of unit tile edges on the boundary of an order-$k$
supertile.  Reading the child multiset off the substitution table --- the eight
placement slots are the same for every species, and slot~$7$ is $\Gamma$ in
every rule, while $\Gamma$ alone has a null slot and therefore seven children
--- and subtracting each glued edge twice gives, with $Q$ the $\Gamma$
perimeter and $G$ the number of edges shared between two \emph{different}
children,
%
\begin{equation}
\label{eq:bdyid}
P(k{+}1)=7P(k)+Q(k)-2G(k{+}1),
\qquad
Q(k{+}1)=6P(k)+Q(k)-2G_{\Gamma}(k{+}1).
\end{equation}
%
The coefficients are read from the rule, not fitted to counts, and~\eqref{eq:bdyid}
is exact at every order computed.  Two structural facts make this small enough
to state.  First, all eight non-$\Gamma$ species have the \emph{same} boundary
sequence: the outline of a supertile does not depend on its type, which is what
collapses a nine-type problem to two.  Second, the glue decomposes by slot pair
into thirteen contact sequences, and every one of them satisfies the same
recurrence as $P$ itself.  The contact set is closed under the operator, which
is why $P$ is.

\begin{fact}
\label{fact:bdyrec}
$P(k)=5P(k-1)-3P(k-2)-P(k-3)$, giving $P=14, 46, 182, 758, 3198, 13534$.  Checked by the Lean kernel,
and the sequence itself counted twice --- see Fact~\ref{fact:bdyexact}.
\end{fact}

\noindent
The seed of that recurrence is counted from placed tiles, and the count used
a rounding key: two edges were identified when their endpoint coordinates
agreed to four decimals.  A key of that kind can fail in both directions, by
merging edges that differ below the tolerance and splitting edges that agree
above it, and nothing inside a floating-point computation bounds the failure
--- it can only be removed.

\begin{fact}
\label{fact:bdyexact}
The same sequence, recounted with edges identified by \emph{equality} of
exact $\mathbb{Q}(\sqrt3)$ endpoint coordinates and from an independent
transcription of the placement chain, agrees at every order computed.  The
rounding key is therefore not load-bearing for any of the boundary results.
\end{fact}

\noindent
The factorisation of $\chi_{\partial}$ is not left as a spectral remark.  Its
quadratic factor applied to the sequence annihilates the growing part and
leaves behind the constant belonging to the marginal eigenvalue, and that
identity is proved for every $k$:

\begin{fact}
\label{fact:bdyinv}
$P(k+2)-4P(k+1)-P(k)=-16$ for all $k$, by induction in Lean.  Together with
the seed this is equivalent to Fact~\ref{fact:bdyrec}, and it is the $+1$
eigenvalue made arithmetic --- the boundary's analogue of $Q_{+}$.
\end{fact}

\begin{fact}
\label{fact:bdychi}
The companion matrix of that recurrence --- the boundary transfer operator ---
has characteristic polynomial $\chi_{\partial}(x)=(x - 1)(x^{2} - 4 x - 1)$.  The marginal eigenvalue
$+1$ is a conserved boundary quantity, of the same kind as the $\pm1$ of
Section~\ref{sec:charges}.
\end{fact}

\begin{fact}
\label{fact:nu}
The Perron root is $\nu=2 + \sqrt{5}=\varphi^{3}=4.23606797750$.
\end{fact}

\noindent
So the area of a supertile inflates by $\lam^{2}=4+\sqrt{15}$ and its perimeter
by $2+\sqrt5$.  This is the $\mathbb{Q}(\sqrt{15})$ / $\mathbb{Q}(\sqrt5)$
divide of Section~\ref{sec:phases} again, but between area and perimeter of the
\emph{same} tile rather than between the Spectre and the Hat.  The supertile
boundary is golden although the tiling is not.

Since $\nu>\lam$ the boundary is fractal: its box dimension is
$\log\nu/\log\lam=1.39925321391$, so a perimeter argument assuming $P(k)\sim\lam^{k}$ is
wrong, and wrong by a definite amount.  Both halves of that comparison are
kernel-checked --- $P(k+3)\ge4P(k+2)+30$ for every $k$, and
$\lam^{2}<16$ hence $\lam<4$ --- but the inference between them is not.
Relating an integer growth bound to a real inflation rate needs an order on
$\mathbb{Z}[\sqrt{15}]$ and a square root, neither of which the Mathlib-free
file has, so what Lean certifies is the pair of bounds and the last step is
elementary arithmetic done by hand.

\subsection*{What the golden appearance is not}

$\varphi^{3}$ next to the Hat's $\varphi^{4}$ invites the conjecture that the
Spectre's boundary combinatorics \emph{are} the Hat's.  They are not, and the
refutation is worth more than the conjecture.

\begin{fact}
\label{fact:goldenabsent}
$x^{2}-4x-1$ divides neither $\chi_{M}$ nor $\chi_{H}$, and a search over
integer matrices built from $M$ and $H$ --- sums, differences, powers,
transposes, adjugates, integer shifts --- finds it in none of them.
\end{fact}

\begin{fact}
\label{fact:lucas}
$\varphi^{n}$ has minimal polynomial $x^{2}-L_{n}x+(-1)^{n}$ with $L_{n}$
Lucas.  So $x^{2}-4x-1$ \emph{is} $\varphi^{3}$ and $x^{2}-7x+1$ \emph{is}
$\varphi^{4}$: once two quantities are both powers of $\varphi$, their minimal
polynomials are forced to be adjacent Lucas quadratics and the resemblance
carries no information beyond $\mathbb{Q}(\sqrt5)$.
\end{fact}

\noindent
The asymmetry is worth recording: the Spectre's own unit is not golden, since
$\chi_{M}$ contributes the trace $8$ and $8$ is not a Lucas number.  The three
operators sit at traces $8$, $7$ and $4$ with norms $+1$, $+1$ and $-1$.

What survives is weaker and true.  The ratio $\mu=\lam^{2}/\nu$ has minimal
polynomial $x^{4} + 32 x^{3} - 46 x^{2} - 32 x + 1$ --- degree four, not two --- so it lies in the compositum
$\mathbb{Q}(\sqrt3,\sqrt5)$ and in neither quadratic field alone.

\begin{prop}
\label{prop:compositum}
The Spectre and the Hat meet in $\mu$ as a compositum, not by either field
containing the other.
\end{prop}


\section{Discussion: a multi-scale reading, and a negative result}
\label{sec:kakeya}

Wang and Zahl~\cite{WangZahl} prove that a family of $\delta$-tubes in
$\mathbb{R}^{3}$ whose directions do not concentrate --- not too many tubes
contained in a common convex set --- has a union of almost maximal volume.
That is a hypothesis about a direction set and a conclusion about a volume, and
both halves are measurable on a finite family.  The braided tiling supplies one:
every shared edge carries a two-strand braid, each strand is a curve in
$\mathbb{R}^{3}$, and its $\delta$-neighbourhood is a tube.  


\subsection*{Directions: the monotile does worse than a lattice}

The obvious bridge fails.  Measured against a periodic hexagonal control with
identical machinery, tube families built from the Spectre are \emph{not} better
spread: on the rank of the dual point cloud the periodic lattice scores higher,
and its slab fractions decay faster.  The reason is not subtle in retrospect.
Kakeya non-concentration asks for a direction set that is \emph{spread};
aperiodicity supplies \emph{non-repetition}.  A hexagonal lattice has a
perfectly spread direction set by symmetry, and a Spectre tiling draws its
edges from a small set of directions on the same hexagonal grid --- repeating
them less often does not make them more numerous.  

\subsection*{Scales: what the substitution actually has}

What a substitution tiling has is scales, and the ladder is the one this paper
is about: a patch of diameter $D$ carries $\rho_{n}=D\lam^{-n}$, graded by the
same unit as everything else.  Along that ladder the families do separate.
Comparing the local direction law inside a $\rho$-ball with the global one, and
normalising against a bootstrap draw of the same size --- without which the
statistic measures the central limit theorem and not the geometry --- the
Spectre holds a ratio near unity, flat across scales spanning a factor of
twenty, while the periodic control sits near a third of that and equally flat.
The Spectre's balls look like random subsamples of the whole patch at every
zoom; the lattice's are hyperuniform, carrying the global direction mix with
sub-random discrepancy.

Flatness across scales is the signature of a \emph{sticky} family, and
stickiness is the obstruction a Kakeya argument must survive rather than a
property that helps it.  So the honest statement is that the Spectre braid is a
concrete finite sticky family, which is a harder object to offer than a better
one.

\subsection*{The exponent the ladder wants}

A genuine multi-scale family puts tubes of many thicknesses in at once, one per
substitution level, with $\delta_{c}=\delta_{0}\lam^{-sc}$.  The natural guess
$s=2$, one factor of the fundamental unit per level, is wrong, and
Section~\ref{sec:bdyop} says why: every interface edge is a \emph{unit} tile
edge whatever level it sits at, so length does not scale with level and only
the population does.  One factor of the unit per level is the right grading for
areas and the wrong grading for a thickness attached to a fixed-length edge.

\begin{prop}
\label{prop:sstar}
Interfaces at coarseness $c$ in a depth-$N$ patch number
$\bigl(T(c)P(N-c)-P(N)\bigr)/2$, so they grow in $c$ by $\mu=\lam^{2}/\nu$, and
the thickness exponent making every level carry comparable volume is
%
\begin{equation}
\label{eq:sstar}
s^{*}=1-\frac{\log\nu}{\log\lam^{2}}=0.300373393044 .
\end{equation}
\end{prop}

\noindent
Two cautions attach to~\eqref{eq:sstar}.  It is exact but \emph{not algebraic}:
a ratio of logarithms of algebraic numbers has a closed form and arbitrary
precision but is not a root of any polynomial.  The algebraic object is $\nu$;
$s^{*}$ is what remains after taking its logarithm, and a reader wanting an
element of $\mathbb{Z}[\sqrt{15}]$ should stop at $\nu$ and accept it in
$\mathbb{Q}(\sqrt5)$.  And $s^{*}$ had to be derived rather than fitted:
because the interface count depends on $T(c)$ and on $P(N-c)$ together, a
single-exponent fit needs $c$ large and $N-c$ large at once, which no finite
patch provides.  Fitting one anyway returns values that drift with depth and
never settle --- and the closed form reproduces those drifting values, which is
how we know the drift was the estimator and not the geometry.


\section{Verification: computed twice, and proved once}
\label{sec:lean}

Computing a Fact twice by code-disjoint routes is a good discipline and it is
still two computations; two programs can be wrong in the same way.  The closed
Facts are therefore also \emph{proved}.  \code{spectrelean.py} emits them as
Lean~4 and the kernel checks the file --- 51 theorems, covering the charge
eigenvector identities, the census totals and their recurrence, the Mystic
identity of Fact~\ref{fact:mystic}, the arithmetic of $\mathbb{Z}[\sqrt{15}]$
including the Pell unit and both channel identities, the frequency
eigenvector, the X-charge parity, and the Seifert minors of
Fact~\ref{fact:sandwich}.

The file is Mathlib-free: everything is \code{Int}, \code{List Int} or a small
structure over them, so it checks against a bare toolchain in seconds.  That
costs some elegance and buys the property that the proof is as auditable as
the arithmetic it certifies.
\begin{footnotesize}
\begin{lstlisting}[language=lean]
theorem q_minus_alternates : vecMul Qminus M = scale (-1) Qminus := by decide

theorem pell_fundamental_unit : lam2 * lam2inv = (1, 0) := by decide

theorem seifert_minor_closed_form (m : Nat) : D m = Dclosed m :=
  (D_closed_pair m).1
\end{lstlisting}
\end{footnotesize}

\noindent
The axiom audit is measured rather than asserted: \code{\#print axioms} is run over every theorem and the result written into the file header. No proof is admitted, and nothing depends on \code{Classical.choice}, so every proof in the file is constructive.  What the kernel does \emph{not} see is the subject of
Section~\ref{sec:boundary}, and is stated in the Lean file as comments rather
than theorems, because a comment cannot be mistaken for a certificate.


\section{The initial surface and the Weyl curvature: a conjecture}
\label{sec:weyl}

This section contains no computation.  It states one conjecture, says what
would support it, and lists what it does not claim, because the question it
touches --- what the substrate has to do with gravity --- is the one a reader
will ask next.

\subsection*{What is known, from outside this paper}

Penrose's Weyl curvature hypothesis asks that the Weyl tensor vanish at the
initial boundary of the universe, so that the gravitational field carries no
free degrees of freedom at the start and the low entropy of the early universe
is the low entropy of the gravitational field~\cite{Penrose1979}.  It is a
hypothesis, not a theorem, and it concerns the Weyl tensor, not the Ricci
tensor: a Weyl-flat initial state still has gravity in the sense of the
Friedmann equations.  In this paper ``no free gravitational field'' is meant,
not ``no gravity''.  That the initial surface is a smooth two-dimensional
carrier and not a point of infinite density is the premise of
Section~\ref{sec:screen}, not something the Weyl hypothesis implies.  The
absence of a time-asymmetric singularity at the start has a precedent in the
Janus-point solutions of the Newtonian $N$-body problem, in which a
scale-invariant measure of structure grows away from a minimum in both time
directions~\cite{BarbourLourenco}; that result is Newtonian and is cited as a
precedent for the shape of the argument, not as support for it.

The sector of general relativity with one chirality of the Weyl tensor
removed is self-dual gravity~\cite{Plebanski1975,Penrose1976}.  A recent
preprint shows that single-minus tree amplitudes of self-dual gravity, which
had been taken to vanish, are nonzero on the half-collinear locus, that is, in
split signature or for complexified momenta, and derives them from a
recursion and, under analyticity assumptions, from an $Lw_{1+\infty}$ Ward
identity~\cite{Guevara2026}.  That preprint is unrefereed.  Two features of it
matter here: it concerns one Weyl chirality, and its kinematics are not
Lorentzian.

\subsection*{The conjecture}

\begin{conj}[Chirality of the initial Weyl data]
\label{conj:weyl}
On the isolated conformal screen the strictly chiral Spectre phase carries one
chirality of the gravitational field, the self-dual sector, and the
Weyl-flat initial state is the state in which the other chirality has not yet
switched on.  The transition of Conjecture~\ref{conj:chirality}, from the
Spectre phase to the reflection-admitting Hat phase, is where the second
chirality appears: parity is restored by embedding in the bulk exactly as
both Weyl chiralities are present in the embedded geometry.
\end{conj}

It is a restatement of Conjecture~\ref{conj:chirality} in the language of the
Weyl spinors: parity-forbidden on the screen corresponds to
one-chirality-only in the gravitational sector.  The existing integer data
fit the shape of this statement and do not test it.  The sign of $Q_{-}$
alternates with the handedness of a patch (Proposition~\ref{prop:parity}), and
the X-charge defines a mirror channel that the strictly chiral screen does not
have (Section~\ref{sec:mass}); the conjecture would say that these are two
faces of one fact about which Weyl chirality is available.


\subsection*{A computable precedent: Barbour's variety on substitution patches}

Barbour and Louren\c{c}o define the \emph{variety} of $N$ point masses, with
the masses normalised to sum to one, as $V=\ell_{\rm rms}/\ell_{\rm mhl}$ with
$\ell_{\rm rms}^{2}=\sum_{i<j}m_im_jr_{ij}^{2}$ and
$\ell_{\rm mhl}^{-1}=\sum_{i<j}m_im_j/r_{ij}$~\cite{BarbourLourenco}.  It is
dimensionless, unchanged by translation, rotation, reflection and rescaling,
and in the zero-energy, zero-momentum Newtonian $N$-body problem it grows away
from the Janus point in both time directions~\cite{BarbourLobo,BarbourEntropy}.
The question put to the substrate is only whether the depth-$n$ Spectre patch,
with one equal mass at each tile centroid, shows anything in this measure that
a control does not.  Table~\ref{tab:variety} gives $V$ for the Spectre patch,
for a hexagonal lattice tuned to give the same number of points in the same
region, and for uniform random points in the same region; the region is the
union of the tiles, so that the outer shape of the non-convex patch is held
fixed and is not what is compared.

\begin{table}[h]
\centering
\caption{Barbour's variety on the depth-$n$ Spectre patch and two controls with
the same region and (up to the lattice rounding) the same $N$.  Random: mean and
standard deviation over 8 draws.  Floating-point computation
(\texttt{spectrevariety.py}), not an exact Fact.}
\label{tab:variety}
\begin{tabular}{rrrrr}
\toprule
depth & $N$ & Spectre & lattice & random \\
\midrule
1 & 9 & 0.3607 & 0.3577 & $0.4923\pm0.0732$ \\
2 & 71 & 0.5387 & 0.5388 & $0.6233\pm0.0205$ \\
3 & 559 & 0.6034 & 0.6017 & $0.6379\pm0.0043$ \\
4 & 4401 & 0.6259 & 0.6254 & $0.6396\pm0.0013$ \\
\bottomrule
\end{tabular}
\end{table}

\textbf{Computation, not Fact.}  The values are floating-point and the
analytic controls (two equal masses, $V=1/8$; an equilateral triangle,
$V=\sqrt3/9$) are the only exact checks.  Three things can be read off.
First, $V$ rises with depth for the Spectre and for both controls alike, from
0.361 at depth 1 to 0.626 at depth 4; that rise is the approach of $N$ points
to a continuum value and is not a property of the substitution.  Second, at
depth 4 the Spectre value and the lattice value differ by 0.09\%, which is
the difference between two nearly equal numbers and shows no distinguishing
signal.  Third, the random control lies 2.2\% above the Spectre, 10
standard deviations of its own spread: in this measure the patch sits at the
lattice-like, uniform end and not at the clustered end.  A different statistic
(the direction law of Section~\ref{sec:kakeya}) was compared there, and the two
are not reconciled here.

What this does not show is as important.  Depth is not time, and a sequence of
patches of increasing $N$ is not one system evolving; Barbour's growth is a
statement about dynamical solutions, so the rise of $V$ with depth is not an
arrow of time and is not evidence for reading the substitution as a Janus-point
flow.  The computation contributes a well-defined, scale-invariant measure
evaluated on the substrate, and the finding that the substrate is
indistinguishable from the lattice in it.


\subsection*{A neighbouring exact structure: the cusp fibre of the complex six-sphere}

One exact structure outside this paper touches the tiles, and it is stated here
so that its limits are on the page.  Alpöge's construction of a complex
structure on the six-sphere~\cite{Alpoge}, a preprint that is unrefereed,
builds a compact complex threefold fibred by tori over $\mathbb{P}^1$ whose cusp
fibre $W$ is a quotient of a toric surface associated with the $A_2$
triangulation of the plane.

\textbf{Imported fact (computed in \texttt{pyCICY}, not recomputed in this
paper, and not one of its numbered Facts).}  Every Tile$(a,b)$, the Hat and the
Spectre among them, is a union of kites of the $[3.4.6.4]$ Laves lattice, which
\texttt{pyCICY.monotile} builds exactly in $\mathbb{Q}(\sqrt3)$.  At the regular
proportions, the $61$ sites of a two-ring Laves patch are all cells of the $A_2$
triangulation of the matching kind, and its $102$ bonds are exactly the $102$
cover relations between them, so the Laves lattice is the Hasse diagram of the
cusp fan.  Under this identification the lattice rolled up on a torus is $W$:
a hexagon centre is a $\mathrm{dP}_6$ component, an edge midpoint a double
curve, a hexagon vertex a triple point.  Per fundamental domain there is one
component, three double curves and two triple points, so $e(W)=2$ counts the
hexagon vertices, and the dual complex is a triangulated torus
($V-E+F=0$).  The check is exact; its limits are that it concerns the
regular proportions $a=\sqrt3\,b$ and not the Spectre's $a=b$, where the
combinatorics survive and the metric match does not.

\textbf{The contrast.}  The substrate of a \emph{periodic} tiling descends to a
torus, which is what $W$ is.  No aperiodic tiling can: it is invariant under no
lattice.  Both objects inherit the lattice from $A_2$; the six-sphere knows
nothing of monotiles, and the relation is a contrast, not an identity of
tilings.

\textbf{What does not connect.}  The period field of the construction is
$\mathbb{Q}(\zeta_{12})$, which shares $\mathbb{Q}(\sqrt3)$ with the tile
geometry and nothing with the field $\mathbb{Q}(\sqrt{15})$ of this paper's
inflation or with the Spectre's $\mathbb{Q}(\sqrt6,\sqrt{10})$.  The integer
$6$ occurs on both sides, in the invariant form of the six-sphere's monodromy
and in the Spectre's inflation $\lambda=(\sqrt6+\sqrt{10})/2$, but it arises
differently in each and nothing uses either; the Laves valences $(6,4,3)$ echo
the orbifold orders $(3,4,\infty)$, but the valence-$4$ sites are double curves
over the cusp and not over the order-$4$ point.  Both are coincidences, and no
link between the two objects should be inferred from them.  The base orbifold
of the six-sphere must be hyperbolic (no Euclidean triangle group contains both
$3$ and $4$), while the substrate's symmetry is Euclidean; and the threefold
carries no holomorphic $3$-form, so it supports no four-dimensional
$\mathcal N=1$ compactification, and nothing here connects it to the screen or
to the Weyl-curvature conjecture above.

\subsection*{What it does not claim, and what is missing}

\begin{itemize}
\item \emph{No map is given.}  There is no construction taking a tile, a
substitution level or a charge to a Weyl spinor or a self-dual curvature.
Without that the conjecture is a statement about analogy of structure.
\item \emph{The signature does not transfer.}  The amplitude results live in
split signature or complexified kinematics.  Nothing here connects them to a
real Lorentzian initial slice, and the self-dual sector is a complexified
subset of Einstein gravity, not a real cosmological state.
\item \emph{No Standard Model content.}  The substrate supplies integer
charges and a tile-frequency spectrum with degeneracy pattern
$1,1,3,2,2$ on nine species; none of this reproduces the field content of the
Standard Model.  The repository's own mass-ratio scan, run with its
false-positive accounting, returns three matches against about three expected
by chance, and is read as numerology.  The substrate is not claimed to
explain any particle.
\item \emph{No cosmological prediction follows.}  The Weyl hypothesis has
observational consequences only through a full cosmological model, and the
substrate does not supply one.
\end{itemize}

\subsection*{What would support or retire it}

The conjecture would be supported by a construction that produces a
one-chirality curvature from the tile data, with the Hat phase supplying the
opposite chirality; and it is retired if the integer chirality data
(Proposition~\ref{prop:parity}) can be shown to carry no relation to the
chirality of any gravitational sector.  Neither has been attempted.


\section{Results: A modelling step}
\label{sec:data}

Isotropic cosmic birefringence is measured nonzero: the 2026 joint
ACT~DR6~$+$~Planck~PR4 analysis \cite{Eskilt2026} gives
$\beta=0.277^{\circ}\pm0.057^{\circ}$.  Exact parity-evenness, the branch that
would disfavour Conjecture~\ref{conj:chirality}, is not what the sky shows ---
but the consistency is generic, since any parity-violating photon coupling
produces it.

The substrate has one intrinsic chiral quantity with a preferred magnitude,
the order-parameter maximum of Fact~\ref{fact:area}.  Identifying it with a
Stokes-level asymmetry and reading the result in degrees gives
\[
\beta_{\rm pred}=2\cdot\frac{10g}{8+10g}\ \mathrm{[deg]}=\frac{70}{67} - \frac{40 \sqrt{15}}{201}=0.27403317\ldots^{\circ},
\]
$0.05205\sigma$ from the measurement.

That agreement is not evidence.  The unit postulate --- that a dimensionless
area fraction is an angle in degrees --- is the entire modelling step and it is
unjustified: degrees are a human convention and nothing in
$\mathbb{Q}(\sqrt{15})$ knows about them.  A closed form landing near a
measurement is what one should expect occasionally from a theory with several
dimensionless constants and a free choice of units.  The substrate supplies a
number, the identification is a guess, and the guess is unrefuted.

It is recorded rather than deleted because it is scheduled: at the
$0.05^{\circ}$ calibration class it separates from zero at $5.5\sigma$, and at
$0.01^{\circ}$ from the WMAP$+$Planck value $0.342^{\circ}$ at $6.8\sigma$.  A
guess with a date on it beats a guess without one and loses to a derivation ---
and Section~\ref{sec:phases} now offers a test that needs no postulate at all.


\section{Conclusion: The boundary, restated}
\label{sec:boundary}

The companion series distinguishes several ways a quantity can fail to be
available; the categories carry different obligations, so they are collected
here rather than left scattered.

\begin{enumerate}
\item \emph{Not a computation at all.}  Conjecture~\ref{conj:chirality} is a
selection statement --- which phase is admissible where, on the grounds that
reflection is ambiently realizable in one setting and not the other.  It is
not a dynamics and does not purport to be.

\item \emph{Exact in principle, not computed here.}  The Spectre's chiral
twist angle, the $\mathbb{Q}(\sqrt{15})$ analogue of Fact~\ref{fact:chiral},
which is why Section~\ref{sec:phases} states it as a prediction.  Likewise the
collared frequency $p$ of Fact~\ref{fact:pdrift}: measurable to any depth,
converging slowly, not yet closed.

\item \emph{Not determined by the data carried.}  The deficit scale $\delta$:
the ratios $\{0,1,2,7\}$ are ours, the scale is not.  And the unit postulate
of Section~\ref{sec:data}.  In both the constraint is exact, the choice is
absent, and what is reported is the constraint.

\item \emph{Obstructed by the transcription.}  That~\eqref{eq:M} transcribes
the published substitution.  Every Fact is downstream of it and none bears on
it.  It is audited --- the geometric census rebuilds the counts from placed
tiles by a route sharing no code, and agrees at every depth --- and an audit is
not a proof.  Section~\ref{sec:bdyop} adds a second input of exactly this kind:
the three seed values of Fact~\ref{fact:bdyrec} are counted from placed tiles,
not derived here.  They are audited the same way --- and, since
Fact~\ref{fact:bdyexact}, audited better than the matrix is, because the
second count is in exact quadratic-field arithmetic rather than a second
float computation.  That removes the rounding key from the boundary results;
it does not remove the transcription, which is the same act of reading a
published substitution and is what this category is for.

\item \emph{Established in the aggregate, open in the fine decomposition.}
The boundary operator of Fact~\ref{fact:bdychi} governs the perimeter, the
glue, and all thirteen contact sequences exactly.  Resolved further, by species
and edge index, a handful of components obey neither it nor its obvious
extensions, while their ratios still approach $\nu$.  So the extra spectrum is
subdominant and cannot move the Perron root, and its exact content is not
settled by the orders we can compute.  What is claimed is the aggregate
operator; the full fine operator is not claimed.

\item \emph{Fixed by reading, not by the geometry.}  The boundary sequences
must be read from $k=1$.  At $k=0$ a supertile is a single tile, there is no
assembly, and a recurrence describing assembly has nothing to describe; read
from $k=0$ only $P$ obeys it, read from $k=1$ all of them do.  Two contact
sequences that look like exceptions differ from their partners in the first
term alone.  This is a convention made explicit rather than a gap, and it is
listed because an unstated one would look like a gap.
\end{enumerate}

\noindent
The fourth category is where the Lean layer does not reach, and the correction
history of Section~\ref{sec:mass} shows that cited results carry it too.  This
is the one place the paper asks to be trusted, and it is better said plainly
than left for a kernel certificate elsewhere to imply otherwise.  A proof about
$M$ is a proof about $M$.


\appendix
\footnotesize

\section*{Declarations}

\subsection*{Funding}
The author received no specific funding for this work.

\subsection*{Author Contributions}
B.H. conceived the theoretical framework, performed the analytical calculations, developed the numerical simulations, analyzed the results, and authored the manuscript.

\subsection*{Competing Interests}
The author declares no competing interests.

\subsection*{Affiliations}
The author declares no affiliations.

\subsection*{Code and Data Availability}
\label{sec:supp}

All source code and data supporting this study, including the Lean~4 sources, are
available at \url{https://github.com/brentharts/spectre/} and
\url{https://github.com/brentharts/CICY/}.


\begin{thebibliography}{99}

\bibitem{SmithEtAl} Smith, D., Myers, J.S., Kaplan, C.S. and
Goodman-Strauss, C. (2024). An aperiodic monotile.
\emph{Combinatorial Theory}, 4(1).

\bibitem{SmithEtAl2} Smith, D., Myers, J.S., Kaplan, C.S. and
Goodman-Strauss, C. (2024). A chiral aperiodic monotile.
\emph{Combinatorial Theory}, 4(2).

\bibitem{bullmore}
E.~Bullmore, M.~Brammer, I.~Harvey, R.~Persaud, R.~Murray, M.~Ron (1994)
Fractal analysis of the boundary between white matter and cerebral cortex in magnetic resonance images: a controlled study of schizophrenic and manic-depressive patients
\url{https://doi.org/10.1017/s0033291700027926}

\bibitem{DorauMuch} P.~Dorau and A.~Much, \emph{From Quantum Relative Entropy to the Semiclassical Einstein Equations}, Phys.\ Rev.\ Lett.\ \textbf{136}, 091602 (2026); arXiv:2510.24491.
\url{https://doi.org/10.1103/lmq8-nsty}




\bibitem{Moritake2026} Moritake, Y., Takiguchi, M., Aihara, T. and
Notomi, M. (2026). Chiral diffraction from aperiodic monotile structure.
\emph{Nature Communications} \textbf{17}, 6085.
\url{https://doi.org/10.1038/s41467-026-75023-7}
\bibitem{BH} M.~Brittenham and S.~Hermiller (2025). Unknotting number is not
additive under connected sum. arXiv:2506.24088v2 (15 September 2025).
\url{https://arxiv.org/abs/2506.24088}
Cited in its corrected version; see~\cite{WZ} and the discussion in
Section~\ref{sec:mass}.
\bibitem{WZ} C.~Wang and Y.~Zhang (2025). A remark on the counterexample to
the unknotting number conjecture. arXiv:2507.14265.
\url{https://arxiv.org/abs/2507.14265}
\bibitem{WangZahl} H.~Wang and J.~Zahl (2025). Volume estimates for unions
of convex sets, and the Kakeya set conjecture in three dimensions.
\emph{arXiv:2502.17655}.

\bibitem{Scharlemann} M.~Scharlemann (1985). Unknotting number one knots are
prime. \emph{Invent. Math.} 82, 37--55.

\bibitem{Penrose1979} R.~Penrose (1979). Singularities and time-asymmetry.
In \emph{General Relativity: An Einstein Centenary Survey}, ed.\ S.~W.~Hawking
and W.~Israel, Cambridge University Press, 581--638.

\bibitem{Plebanski1975} J.~F.~Pleba\'nski (1975). Some solutions of complex
Einstein equations. \emph{J. Math. Phys.} 16, 2395--2402.

\bibitem{Penrose1976} R.~Penrose (1976). Nonlinear gravitons and curved
twistor theory. \emph{Gen. Relativ. Gravit.} 7, 31--52.

\bibitem{Guevara2026} A.~Guevara, A.~Lupsasca, D.~Skinner, A.~Strominger and
K.~Weil (2026). Single-minus graviton tree amplitudes are nonzero. OpenAI
preprint (unrefereed). \url{https://cdn.openai.com/pdf/graviton.pdf}

\bibitem{Alpoge} L.~Alp\"oge (2026). A compact complex threefold fibred by
tori over the projective line, and the six-sphere. Preprint (unrefereed).
\url{https://alpo.ge/s6.pdf}

\bibitem{BarbourLobo} J.~Barbour, F.~S.~N.~Lobo and M.~I.~R.~Louren\c{c}o
(2026). Structural morphology and the gravitational arrow of time.
arXiv:2607.27526.

\bibitem{BarbourEntropy} J.~Barbour (2021). Entropy and cosmological arrows of
time. arXiv:2108.10074.

\bibitem{BarbourLourenco} J.~Barbour and M.~I.~R.~Louren\c{c}o (2026). A
scale-invariant theory of the universe. arXiv:2608.05929.

\bibitem{Eskilt2026} Eskilt, J.R. \emph{et al.} (2026). Joint ACT DR6 and
Planck PR4 constraints on cosmic birefringence.

\end{thebibliography}

\end{document}
